\newglossaryentry{Bias}{
    name={Bias (Sesgo):},
    text={bias},
    description={Presencia de distorsiones sistemáticas en los datos, algoritmos o procesos de un sistema de inteligencia artificial, que pueden influir en el aprendizaje y conducir a resultados injustos o desbalanceados entre distintos individuos o grupos. Estos sesgos pueden originarse en múltiples etapas del ciclo de vida del sistema, incluyendo la recolección de datos, el diseño del modelo y la interacción con los usuarios}
}

\newglossaryentry{Fairness}{
    name={Fairness (Equidad):},
    text={fairness},
    description={Propiedad de un sistema de inteligencia artificial que implica la ausencia de prejuicio o favoritismo hacia individuos o grupos en el proceso de toma de decisiones, garantizando resultados justos y no discriminatorios, especialmente respecto a atributos sensibles}
}

\newglossaryentry{Generalización}{
    name={Generalización:},
    text={generalización},
    description={Capacidad de un modelo de mantener un buen desempeño al ser aplicado a datos nuevos o no vistos previamente, especialmente en contextos distintos a aquellos utilizados durante su entrenamiento}
}

\newglossaryentry{Modelo Fundacional}{
    name={Modelo Fundacional:},
    text={modelo fundacional},
    description={Modelo de inteligencia artificial entrenado sobre grandes volúmenes de datos, generalmente de forma auto-supervisada, capaz de adaptarse a múltiples tareas específicas mediante técnicas como \textit{fine-tuning} o \textit{zero-shot learning}.}
}

\newglossaryentry{CT}{
    name={Tomografía Computarizada (Computed Tomography, CT):},
    text={tomografías computarizada},
    description={Técnica de imagen médica que utiliza rayos X para obtener imágenes transversales detalladas del cuerpo, permitiendo visualizar estructuras internas con alta resolución y apoyar el diagnóstico clínico}
}

\newglossaryentry{WSI}{
    name={Whole Slide Image (WSI):},
    text={Whole Slide Image},
    description={Imagen digital de alta resolución obtenida a partir de muestras histopatológicas completas, que permite el análisis computacional de tejidos a nivel microscópico en patologías como el cáncer}
}

\newglossaryentry{IA}{
    name={Inteligencia Artificial (IA):},
    text={inteligencia artificial},
    description={Disciplina de la informática que desarrolla sistemas capaces de realizar tareas que normalmente requieren inteligencia humana, como el aprendizaje, el razonamiento y la toma de decisiones}
}

\newglossaryentry{AA}{
    name={Aprendizaje Automático:},
    text={aprendizaje automático},
    description={Subcampo de la inteligencia artificial que se enfoca en el desarrollo de algoritmos que permiten a los sistemas aprender patrones a partir de datos y mejorar su desempeño sin ser programados explícitamente}
}

\newglossaryentry{DL}{
    name={Deep Learning:},
    text={deep learning},
    description={Enfoque dentro del aprendizaje automático basado en redes neuronales profundas, capaz de modelar relaciones complejas en grandes volúmenes de datos}
}

